\documentclass[11pt]{article}
\topmargin -.25in
\oddsidemargin 0in
\textwidth 6.5in
\textheight 9in
\def\v{\vspace{11pt}}
\def\d{\displaystyle}
\pagestyle{empty}
\usepackage{graphics,amsmath,amssymb,ulem,tikz}

\begin{document}
\centerline {\bf \large Quiz \# 3}
\v

\noindent Math 118 \hfill Names: \rule{2in}{.25mm}

\noindent Instructor: \rule{2in}{.25mm} \hfill \rule{2in}{.25mm}

\noindent Fall 2025 \hfill \rule{2in}{.25mm}

\hfill \rule{2in}{.25mm}

\v

\noindent To receive full credit you must show all work and provide clear, complete explanations.  Use of books, notes, calculators, etc.\ is not permitted.  

\begin{enumerate}
\item (5 points) Find the center and radius of the circle:
    \begin{equation*}
        x^2 + 4x + y^2 - 6y + 5 = 0
    \end{equation*}


\newpage

\item (5 points) 
\vspace{0.1in}
\begin{enumerate}
    \item In slope intercept form, write the equation of the line that passes through the points $(-2,2)$ and $(7,20)$. \vspace{3.5in}
    \item Write the equation of the line that is perpendicular to AND has the same $y$-intercept as the line that you computed in part (a).
\end{enumerate}

\vspace{4in}

\item (5 points) Find the domain of the function:
\begin{equation*}
    f(x) =  \frac{x}{\sqrt{x^2 - 4x-21}}
\end{equation*}

\newpage

\item (5 points) Recall that the difference quotient of a function $f(x)$ is:
\begin{equation*}
    \frac{f(x+h) - f(x)}{h}
\end{equation*}
Find the difference quotient of $\displaystyle f(x) = \frac{1}{x}$.
\end{enumerate}
\end{document}